\hypertarget{advanced-concurrency-shared-memory-and-proxies}{%
\chapter{Advanced Concurrency: Shared Memory and
Proxies}\label{advanced-concurrency-shared-memory-and-proxies}}

Building on Chapter 27, this chapter explores the high-performance
communication mechanisms required for massive scale data processing in
Python.

\hypertarget{multiprocessing.shared_memory-zero-copy-communication}{%
\subsection{\texorpdfstring{66.1
\texttt{multiprocessing.shared\_memory}: Zero-Copy
Communication}{66.1 multiprocessing.shared\_memory: Zero-Copy Communication}}\label{multiprocessing.shared_memory-zero-copy-communication}}

Prior to Python 3.8, \passthrough{\lstinline!multiprocessing!} relied on
pickling objects and sending them via pipes/sockets, which was slow for
large arrays. \passthrough{\lstinline!shared\_memory!} provides a way to
allocate raw memory that can be accessed by multiple processes without
copying.

\hypertarget{the-sharedmemory-object}{%
\subsubsection{\texorpdfstring{1. The \texttt{SharedMemory}
Object}{1. The SharedMemory Object}}\label{the-sharedmemory-object}}

\begin{itemize}
\tightlist
\item
  \textbf{Creation}: One process creates the memory block with a unique
  name.
\item
  \textbf{Attachment}: Other processes ``attach'' to the memory using
  the name.
\item
  \textbf{Internals}: On POSIX, this uses
  \passthrough{\lstinline!shm\_open()!} and
  \passthrough{\lstinline!mmap()!}. On Windows, it uses
  \passthrough{\lstinline!CreateFileMapping()!}.
\end{itemize}

\hypertarget{shareablelist-and-ndarray-integration}{%
\subsubsection{\texorpdfstring{2. \texttt{ShareableList} and
\texttt{ndarray}
Integration}{2. ShareableList and ndarray Integration}}\label{shareablelist-and-ndarray-integration}}

You can wrap a \passthrough{\lstinline!SharedMemory!} block in a
\passthrough{\lstinline!ShareableList!} (for basic types) or use it as
the buffer for a NumPy array:

\begin{lstlisting}[language=Python]
from multiprocessing import shared_memory
import numpy as np

# Creator
shm = shared_memory.SharedMemory(create=True, size=1024)
arr = np.ndarray((128,), dtype=np.int64, buffer=shm.buf)
arr[:] = np.arange(128)

# Consumer (in another process)
existing_shm = shared_memory.SharedMemory(name=shm.name)
arr_copy = np.ndarray((128,), dtype=np.int64, buffer=existing_shm.buf)
print(arr_copy[10]) # Output: 10
\end{lstlisting}

\hypertarget{managers-and-proxies-distributed-objects}{%
\subsection{66.2 Managers and Proxies: Distributed
Objects}\label{managers-and-proxies-distributed-objects}}

The \passthrough{\lstinline!multiprocessing.Manager!} allows you to
share complex Python objects (like dicts or custom classes) across
processes using a server-client architecture. * \textbf{The Server
Process}: A hidden process manages the ``real'' objects. *
\textbf{Proxies}: Worker processes receive ``Proxy'' objects that look
like the real thing but send every method call over a socket to the
server process. * \textbf{Performance Note}: While flexible, proxies are
much slower than shared memory because every access involves a
network/IPC round-trip and synchronization.

\begin{center}\rule{0.5\linewidth}{0.5pt}\end{center}
